

\section{Généralités sur MapReduce}

\begin{frame}
\frametitle{Qu'est-ce que MapReduce}

C'est un \emph{framework} d'exécution.

\begin{itemize}
  \item Processus en deux phases (devinez?)  appliqué à une collection d'items.
  \item À chaque phase, on applique une fonction fournie par le développeur.
\end{itemize}

\vfill

Quel intérêt?

\begin{itemize}
  \item Implante de manière \emph{générique} une exécution de type "group-by" à la SQL.
  \item \alert{Scalable car facilement exécutable en parallèle}. 
\end{itemize}

\vfill

\alert{Ce n'est pas l'idéal\,!} Très contraint, très peu expressif, lent.

\vfill

Cas particulier d'une \alert{chaîne de traitement} scalable. Principe très ancien
(prog. fonctionnelle), remis au goût du jour par Google (2004) pour traitements massivement distribués. 

\vfill

Les systèmes modernes ont étendu le principe: nous en reparlerons.

\end{frame}


\section{MapReduce, c'est simple}



\begin{frame}[fragile]
\frametitle{Faisons du jus de pomme avec MapReduce}

Vous savez comment on fait du jus de pomme? Voici comment
on fait  avec MapReduce (sans le dire).

\figSlideWithSize{juspommeMR}{10cm}

\begin{defblock}{À retenir}

\begin{itemize}
  \item \alert{L'atelier de préparation}: on transforme les items en entrée, un par un.
        \alert{Traitement individuel, ordre indifférent.}
        
  \item \alert{L'atelier d'assemblage}: on a \alert{groupé} les items, on leur applique
   un \alert{traitement collectif}.
\end{itemize}
\end{defblock}
\end{frame}



\begin{frame}[fragile]
\frametitle{Une petite extension}

Rien n'impose d'avoir du 1 pour 1 dans l'atelier de préparation. On peut
rejeter des pommes pourries, et couper les autres en quartiers.

\figSlideWithSize{juspommeMR+}{11cm}

\begin{defblock}{À retenir}
L'atelier de préparation n'est pas limité à une transformation un pour un des produits consommés.
   Il peut prendre en entrée des produits d'une certaine nature, et sortir
   des produits d'une autre nature.
\end{defblock}
\end{frame}


\begin{frame}[fragile]
\frametitle{Passons au BigData}

Si on veut faire beaucoup de jus de pomme: on parallélise.


\figSlideWithSize{BigjuspommeMR}{9cm}

\end{frame}



\begin{frame}[fragile]
\frametitle{Des remarques importantes}

\alert{Parallélisation = indépendance des données et des transformations}

\medskip
Si l'ordre d'épluchage était important, ce ne serait pas si simple;

\smallskip

Idem si l'épluchage d'une pomme dépendait de l'épluchage des autres. 

\vfill

\alert{Qu'est-ce qu'on gagne?}

\medskip

$n$ fois plus de ressources (cuisinier, matériel),  $n$ fois plus de jus de pomme !

\medskip

\begin{defblock}{Scalabilité}
La production de jus de pomme est parallélisable et proportionnelle aux ressources
   (humaines et matérielles) affectées.
\end{defblock}

\smallskip

Mon processus est robuste!

\begin{defblock}{Robustesse}
Une panne affectant la production de jus de pomme n'a qu'un effet \alert{local}
   et ne remet pas en cause \alert{l'ensemble} de la production.
\end{defblock}

\end{frame}



\begin{frame}[fragile]
\frametitle{Encore mieux: je peux faire des jus de fruit}

Jus de pomme, d'orange ou d'ananas: même procédé + un \alert{tri}.

\figSlideWithSize{jusfruits}{9cm}

\end{frame}



\begin{frame}[fragile]
\frametitle{Encore des remarques importantes}

\alert{Je dois trier et regrouper dans l'atelier d'assemblage}

\medskip

Ma transformation s'applique à des groupes distincts: ils sont
constitués quand on initialise la phase d'assemblage.

\vfill


\alert{Pour être capable de trier, je dois étiqueter les items}

\medskip

C'est le rôle de l'atelier de préparation: on met une étiquette
sur chaque item avant de le transmettre à l'assemblage.


\end{frame}



\begin{frame}[fragile]
\frametitle{Bouquet final: je parallélise aussi l'assemblage}

Et voici du Map Reduce complet.

\figSlideWithSize{jusfruits+}{9cm}

Maintenant, je dois aussi \alert{distribuer} les items: pommes et oranges en
haut, ananas en bas.

\end{frame}


\section{MapReduce, le modèle}


\begin{frame}[fragile]
\frametitle{Résumé (version basique)}

La phase de \texttt{map} (préparation). Elle prend en paramètre une fonction $F_{map}()$.

\begin{itemize}
  \item Le \textit{framework} parcourt la collection en entrée, document par document.
  \item Le \textit{framework} applique à chaque document $d_i$ la fonction  $F_{map}()$.
  \item Le \textit{framework} obtient une valeur $v_i$ qu'il place dans un \alert{accumulateur} $A$.
\end{itemize}

\medskip

La phase de \texttt{reduce} (assemblage). Elle prend en paramètre une fonction $F_{red}()$.


\begin{itemize}
  \item Le \textit{framework} applique $F_{red}()$ aux valeurs contenues  dans l'accumulateur $A$.
  \item Le \textit{framework} obtient une valeur qui est le résultat du processus. 
\end{itemize}

\begin{defblock}{Essentiel}
Les fonctions $F_{map}()$ et $F_{red}()$ sont \alert{la partie "métier", fonctionnelle}. Le reste
est une \alert{infrastructure générique} prise en charge par un \alert{framework}.
\end{defblock}

Dans notre schéma: le cuisinier est $F_{map}()$, le pressoir est $F_{red}()$ 
\end{frame}


\begin{frame}[fragile]
\frametitle{Notions et vocabulaire}

À retenir:

\begin{itemize}
  \item \alert{Item (entrée)}, \alert{document} en ce qui nous concerne.
  
  \item \alert{Fonction de Map}, implante la transformation appliquée à un item pendant la phase de Map.

  \item \alert{Paire intermédiaire}, c'est le résultat de la fonction de Map. 
  
  Une paire intermédiaire $(k, v)$ comprend l'identifiant (étiquette) du groupe auquel
  $v$ appartient.

  \item \alert{Groupe intermédiaire}, ensemble des valeurs produites par la fonction
  de Map et partageant le même identifiant de groupe.
  
    \item \alert{Fonction de Reduce}, implante la transformation appliquée à un groupe
    pendant la phase de Reduce.
    
\end{itemize}

\alert{Ne pas retenir (pour l'instant)}: comment tout cela s'exécute en distribué, avec
parallélisation et reprise sur panne.

\end{frame}


\begin{frame}[fragile]
\frametitle{Vision d'ensemble (en centralisé)}


\figSlideWithSize{mapreduce-mongo}{11cm}

\end{frame}



\begin{frame}[fragile]
\frametitle{Conception d'un processus Map/Reduce: les questions}

Questions pour la phase de Map:
\begin{itemize}
  \item  Documents en entrée. D'où viennent-ils,
    quelle est leur structure? Un traitement MapReduce s'applique à un flux de documents, on
    ne peut rien faire si on ne sait pas en quoi ils consistent.

  \item Groupes ? Combien y en a-t-il? Comment les
    identifier (valeur de clé identifiant un groupe, l'étiquette)? Comment déterminer 
    le ou les étiquettes  dans un document en entrée?

  \item Quelles sont les valeurs intermédiaires que je veux produire à partir d'un document et
    placer dans des groupes? Comment les produire à partir d'un document?
\end{itemize}

Pour la fonction de Reduce, c'est encore plus simple:
\begin{itemize}
  \item Quelle est la nature de l'agrégation qui va prendre un groupe et produire une valeur finale?
\end{itemize}
\end{frame}



\begin{frame}[fragile]
\frametitle{Conception d'un processus Map/Reduce: les fonctions}


Un traitement MapReduce se spécifie sous la forme de deux fonctions

\begin{itemize}

  \item La fonction de Map prend \alert{toujours} en entrée \alert{un} (un seul) document;
   elle produit \alert{toujours} des paires $(k, v)$ où $k$ est l'étiquette (la clé) du groupe
   et $v$ la valeur intermédiaire.
 \item La fonction de Reduce prend \alert{toujours} en entrée une paire $(k, list(v))$,
   où $k$ est l'étiquette (la clé) du groupe
   et $list(v)$ la liste des valeurs du groupe.
\end{itemize}
   
\vfill

À vous de jouer! Exercices.

\end{frame}